\documentclass[10pt]{article}
\usepackage[T1]{fontenc}
\usepackage[utf8]{inputenc}
\usepackage{mathpazo}
\usepackage[scaled=0.9]{helvet}
\usepackage[letterpaper,margin=0.42in]{geometry}
\usepackage{enumitem}
\usepackage{titlesec}
\usepackage{xcolor}
\usepackage{microtype}
\usepackage[colorlinks=true,linkcolor=black,urlcolor=black,citecolor=black]{hyperref}

\titleformat{\section}{\normalsize\bfseries\scshape}{}{0em}{}[\vspace{-4pt}\rule{\linewidth}{0.8pt}]
\titlespacing{\section}{0pt}{4pt}{1pt}
\setlist[itemize]{leftmargin=1.2em,itemsep=0.4pt,topsep=1pt,parsep=0pt,label=\textbullet}
\setlength{\parindent}{0pt}
\pagestyle{empty}
\linespread{0.98}

\newcommand{\role}[2]{\noindent\textbf{#1}\hfill\textbf{#2}\par}
\newcommand{\sub}[1]{\noindent\textit{#1}\par\vspace{1.5pt}}

\begin{document}

\begin{center}
{\LARGE\bfseries Luiz Takahashi}\\[1pt]
{\small l.takahashidosreis@wsu.edu $\mid$ (509) 338-8073 $\mid$ Pullman, WA $\mid$ \href{https://github.com/LTTakahashi}{github.com/LTTakahashi} $\mid$ \href{https://luiztakahashi.com}{luiztakahashi.com}}
\end{center}
\vspace{-3pt}

\section{Education}
\role{Washington State University}{Pullman, WA}
\role{B.S.\ in Computer Science, GPA: 3.84/4.0}{Expected May 2027}
President's Honor Roll (4 semesters). Coursework: Machine Learning, Data Mining, Algorithms, Linear Algebra.

\section{Research Experience}

\role{Independent Researcher, Computational Biology (single-cell / neural organoids)}{2025 -- Present}
\sub{Self-directed; sole-author preprint (Zenodo, 2026).}
\begin{itemize}
\item Conducted a single-author study of the off-target compartment of human neural organoids: mapped a 1.77M-cell organoid atlas onto a harmonized 2.6M-cell fetal-brain reference (scVI / scANVI / scArches), defining calibrated off-manifold and label-entropy scores that resolve a six-state taxonomy of off-target cells.
\item Established an \textit{identifiability boundary}: proved a standard contrastive disentanglement model cannot separate lineage from in-vitro state (an adversary recovers cell source at $>$0.998 accuracy regardless of latent capacity), then recovered confound-invariant geometry via sliced-Wasserstein cancellation, validated with a confound-controlled cross-atlas test.
\end{itemize}

\role{Undergraduate Research Assistant, CASAS Smart Home Lab \& Neuropsychology and Aging Lab}{Aug 2025 -- Present}
\sub{Washington State University. Advisors: Dr.\ Diane J.\ Cook \& Dr.\ Maureen Schmitter-Edgecombe. Funded by NIH/NIA.}
\begin{itemize}
\item Developed a room-aware Kalman filter for indoor localization from sparse smart-home sensors monitoring 35 older adults (ages 60--90+, 4 days to 22 months); reduced localization error by 18\% with 92\% statistical consistency (vs.\ 65\% baseline) across 2M+ sensor events.
\item Found cognitive accuracy significantly associated with behavioral variety ($r=-0.43$, $p=0.015$), driven almost entirely by MCI participants ($r=-0.93$, $p=0.002$); showed the effect dissolves under a noise-multiplied outcome, isolating where the signal lives.
\item Built a LASSO classifier (healthy vs.\ impaired, AUC $=0.78$) identifying spatial entropy ($\beta=-0.42$) and night-activity fraction ($\beta=+0.38$) as clinically interpretable biomarkers consistent with sundowning.
\end{itemize}

\role{Independent Researcher (Columbia University / Wen Lab collaboration), Multi-View Interaction Hierarchy}{2025 -- Present}
\begin{itemize}
\item Developed a hierarchy for multi-view interaction detection with completeness guarantees; proved that the common view-disagreement heuristic cannot detect interactions, and bounded the interaction order recoverable at a given sample size via a Hoeffding functional decomposition.
\item Demonstrated on a 30,000-person GNPC proteomic dataset (9 organ biological-age gaps, 18 disease outcomes) that a seemingly significant cross-organ interaction effect is a between-cohort batch artifact, reproducing the Simpson's paradox from first principles.
\end{itemize}

\role{Research Assistant, Center for Artificial Intelligence (C4AI), University of S\~ao Paulo}{May -- Aug 2025}
\sub{HarpIA Project. Advisor: Dr.\ Fabio Cozman.}
\begin{itemize}
\item Developed SQL Component Match, a semantic natural-language-to-SQL evaluation metric using AST analysis ($\sim$4,000 lines of Python): 95\% semantic accuracy vs.\ a 70\% Spider baseline with $O(n^{1.08})$ scalability; built a 5,686-line test suite (234 methods) with multi-stage normalization, a DNF circuit breaker, and Unicode homograph detection.
\end{itemize}

\section{Publications \& Presentations}
\begin{itemize}
\item \textbf{L.\ Takahashi.} \textit{NullState: a reference-anchored framework for detecting, diagnosing, and quantifying the off-target compartment of human neural organoids.} Preprint, Zenodo, 2026. \href{https://doi.org/10.5281/zenodo.21459720}{doi:10.5281/zenodo.21459720}. (Sole author.)
\item \textbf{L.\ Takahashi}, J.\ Wen. Multi-view interaction hierarchy for health screening. \textit{In preparation} (target: AISTATS 2027).
\item \textbf{L.\ Takahashi}, D.\ J.\ Cook, M.\ Schmitter-Edgecombe. Day-similarity structure in older adults. \textit{In preparation} (target: journal).
\item \textbf{L.\ Takahashi}, D.\ J.\ Cook, M.\ Schmitter-Edgecombe. Room-Aware Behavioral Profiling for Cognitive Health Monitoring in Smart Homes. \textit{Gerontechnology} 25(Suppl.), 2026 (ISG World Conference, Vancouver). \href{https://doi.org/10.4017/gt.2026.25.2.1421.3}{doi:10.4017/gt.2026.25.2.1421.3}.
\end{itemize}

\section{Teaching}
\role{Teaching Assistant, Washington State University}{2025 -- Present}
\begin{itemize}
\item \textbf{Program Design \& Development (C/C++):} lead a weekly 3-hour lab for 15+ students; mentor on data structures, algorithms, and debugging during office hours.
\item \textbf{Data Mining:} sole teaching assistant for the course; graded assignments and exams for 50+ students.
\end{itemize}

\section{Technical Skills}
\noindent\textbf{Languages:} Python (advanced), C (proficient), C++, SQL, C\#, Bash \quad\textbf{ML/AI:} PyTorch, scikit-learn, scvi-tools/Scanpy, Kalman filtering, LASSO/Ridge, KNN, neural networks \\[1pt]
\textbf{Data:} Pandas, NumPy, SciPy, Matplotlib, Seaborn, sqlglot, time-series analysis, synthetic data generation \quad\textbf{Tools:} Git, Docker, Linux, \LaTeX, Jupyter

\end{document}
